\section{从晶体管老化到现场失效}

制造测试回答“这颗器件现在是否符合规格”，RAS 回答“运行中出现错误后怎样检测和恢复”，两者之间还缺少一条时间轴：电场、电流密度、温度和机械应力会在多年运行中逐渐改变晶体管、互连、封装与焊点。可靠性设计的目标不是宣称器件永不变化，而是把工作条件限制在经过验证的寿命模型内，并在裕量耗尽前通过监测、降额或维护阻止错误传播。\cite{nasa-pof-reliability}

\topic{电迁移首先改变的是金属，不是逻辑功能}

\term{电迁移}{electromigration, EM}是导电电子与金属原子交换动量造成的长期物质迁移。在高电流密度连线或过孔中，一端可能形成空洞并使电阻上升，另一端可能堆积形成突起和短路。逻辑仿真仍会把这条网当作理想导线；真实芯片却可能先表现为压降增大、边沿变慢和偶发时序错误，最后才成为永久开路。

常用 Black 模型把平均失效时间的趋势写成
\[
  \mathrm{MTTF}_{EM} \propto J^{-n}
  \exp\left(\frac{E_a}{kT}\right),
\]
其中 $J$ 是电流密度，$n$ 与材料和失效模式有关，$E_a$ 是激活能，$T$ 是绝对温度。该式表达两个工程事实：提高电流密度会缩短寿命，升高温度会加快扩散。公式中的参数必须来自目标工艺与结构的测试，不能用一组通用常数给任意芯片预测精确寿命。

物理设计因此对电源网、时钟主干和高翻转率信号设置电流密度上限，必要时加宽金属、增加并联过孔或分散负载。运行期若持续提高电压和频率，动态电流与结温同时上升，EM 寿命不是只按性能提升比例下降；电流和温度会共同放大损耗。

\topic{BTI、HCI 与 TDDB 消耗晶体管裕量}

\term{偏置温度不稳定性}{Bias Temperature Instability, BTI}在栅极长期偏置和高温下改变界面陷阱与电荷状态，使阈值电压漂移。结果可能是晶体管驱动能力下降、路径延迟增大或 SRAM 读写裕量缩小。部分漂移在撤去应力后可以恢复，长期成分则会积累，因此实际损耗与占空比、温度、电压和时间共同相关。

\term{热载流子注入}{Hot Carrier Injection, HCI}发生在载流子被强电场加速并损伤栅介质或界面时，常见结果同样是跨导与阈值变化。BTI 更关注长期偏置与温度，HCI 更敏感于开关和局部高场；两者都可能先以“最慢路径变得更慢”出现，而不是立即产生一位稳定的错误。

\term{时间相关介质击穿}{Time-Dependent Dielectric Breakdown, TDDB}描述绝缘介质在长期电场下逐步产生缺陷，最终形成泄漏通道或击穿。提高供电电压可以暂时增加速度，却会提高介质电场并加快老化；因此超频稳定测试只能证明短时间内没有越过功能边界，不能证明多年寿命仍满足目标。

\begin{longtable}{L{0.17\linewidth}L{0.23\linewidth}L{0.26\linewidth}L{0.24\linewidth}}
\toprule
老化机制 & 主要应力 & 早期可见症状 & 设计与运行期措施 \\
\midrule
EM & 金属/过孔电流密度、温度 & 电阻上升、IR drop 增大、间歇时序错 & 加宽金属、并联过孔、限流、降温 \\\tablerowrule
BTI & 栅偏置、温度、时间与占空比 & 阈值漂移、路径与 SRAM 裕量下降 & 老化角 STA、时序余量、DVFS 降额 \\\tablerowrule
HCI & 局部高电场与开关活动 & 跨导下降、关键路径逐渐变慢 & 限制电压、器件选型、活动率与布局约束 \\\tablerowrule
TDDB & 介质电场、温度和累积时间 & 泄漏增加，最终短路或功能失效 & 电压额定值、介质规则、保护与寿命测试 \\\tablerowrule
热循环疲劳 & 温度摆幅、循环次数、材料膨胀失配 & 焊点裂纹、接触电阻间歇变化 & 控制温度摆幅、封装/焊点设计、循环试验 \\
\bottomrule
\end{longtable}

\topic{封装与焊点承受的是温度差，而不只是最高温度}

硅、封装基板、焊料和 PCB 的热膨胀系数不同。机器反复从冷机升到满载温度再冷却，各层希望伸缩不同的长度，应力便集中在凸点、焊球、过孔和界面。单次循环造成的塑性应变很小，重复数千次后仍可能形成裂纹。便携设备的开关机、服务器的周期性高负载和风扇控制引起的温度摆动，都会参与这本机械寿命账。

最高结温决定扩散和介质老化速度，温度摆幅与循环次数决定疲劳。把风扇曲线调成在阈值附近频繁升降，平均温度可能不高，却增加热循环；让器件一直维持较高温度又会加快 EM、BTI 和 TDDB。平台控制器需要同时考虑峰值、平均值、变化速率和循环次数，而不是只监控一次是否触发 THERMTRIP。

\begin{pdfdiagram}[x=1cm,y=1cm]
  \node[hw state, minimum width=2.7cm, align=center] (stress) at (0.8,3.45)
    {任务与环境\\电压、电流、温度\\活动率、温度循环};
  \node[hw compute, minimum width=2.7cm, align=center] (physics) at (4.2,3.45)
    {物理损耗\\EM / BTI / HCI\\TDDB / 焊点疲劳};
  \node[hw state, minimum width=2.7cm, align=center] (margin) at (7.6,3.45)
    {裕量变化\\电阻、泄漏、延迟\\接触与噪声余量};
  \node[hw control, minimum width=2.7cm, align=center] (symptom) at (11.0,3.45)
    {可见证据\\纠错率、时序错\\传感器与间歇断连};
  \draw[hw bus] (stress) -- (physics);
  \draw[hw bus] (physics) -- (margin);
  \draw[hw bus] (margin) -- (symptom);

  \node[hw control, minimum width=3.1cm, align=center] (policy) at (7.6,0.85)
    {平台策略\\降压降频、限功率\\迁移负载、维护更换};
  \node[hw state, minimum width=3.1cm, align=center] (evidence) at (4.2,0.85)
    {趋势账本\\位置、应力时间\\错误计数、校准漂移};
  \draw[hw return] (symptom.south) -- (11.0,0.85) -- (policy.east);
  \draw[hw return] (policy) -- (evidence);
  \draw[hw return] (evidence.west) -- (0.8,0.85) -- (stress.south);
\end{pdfdiagram}
\diagramnote{老化不是一次离散故障，而是“应力积累—物理损耗—裕量下降—症状出现”的慢通路。平台只能观测传感器、错误计数和校准漂移等代理量，再通过降额或维护减少后续应力；恢复一次错误不会把已经消耗的材料寿命恢复到初始状态。}

\section{加速寿命试验、降额与现场闭环}

若一颗器件目标寿命为十年，不能等待十年后再决定设计是否合格。\term{加速寿命试验}{accelerated life test}在更高温度、电压、电流密度或更大温度摆幅下施加已知应力，观察失效分布，再用经过验证的物理模型外推到额定工作条件。不同失效机制的加速变量不同：EM 依赖电流密度和温度，TDDB/BTI 强烈依赖电场与温度，封装疲劳依赖温度摆幅和循环。把所有试验都概括成“高温烘烤”会混淆要验证的机制。

外推必须保持失效模式一致。应力过高可能触发额定条件下不会出现的新机制，例如封装材料软化、保护二极管导通或局部热失控；此时得到的是另一种失效，而不是正常使用寿命的加速版本。试验报告至少要保存样本数、删失样本、应力条件、失效判据、物理分析结果和模型参数置信区间。

\topic{浴盆曲线区分缺陷、随机错误和磨损}

早期失效常来自制造缺陷、污染、装联空洞或边缘工艺，筛选和 burn-in 可以提前暴露其中一部分。稳定使用期的失效率相对平缓，软错误、环境扰动与随机部件失效更突出。磨损期则由 EM、介质老化和机械疲劳等累计机制主导，失效率随时间上升。FIT 或 MTBF 常用来描述稳定使用期的统计风险，不能把 MTBF 解释为单台机器的预定寿命，也不能用常数失效率掩盖磨损期。

\topic{从芯片 sign-off 到平台维护要共享同一本应力账}

芯片实现阶段以老化模型做 timing sign-off，并为电源网执行 EM/IR 检查；封装和主板阶段做温度循环、振动、湿热与连接可靠性试验；平台阶段用电压、电流、结温、风扇、ECC 纠错率、链路重训次数和传感器漂移估计实际应力。三层使用的指标不同，却应能回答同一个问题：现场工作条件是否仍落在被验证的寿命包络内。

一个可执行的维护规则不应写成“温度高时关注设备”，而应给出对象、窗口和动作。例如：某 DIMM 在 24 小时内同一器件的可纠正错误持续上升，先把页面迁出并下线对应 rank，保存 syndrome 与温度历史，再安排更换；某 PCIe 插槽的重训次数随热循环增加，则同时检查连接器、焊点、参考时钟和供电，而不是只更换端点卡。

\begin{longtable}{L{0.18\linewidth}L{0.27\linewidth}L{0.27\linewidth}L{0.18\linewidth}}
\toprule
阶段 & 建立的证据 & 决策边界 & 失败后的闭环 \\
\midrule
工艺与器件鉴定 & 试验结构、加速模型、参数分布 & 工艺设计规则与额定条件 & 更新模型、材料或工艺窗口 \\\tablerowrule
芯片与封装 sign-off & 老化角时序、EM/IR、热与机械仿真 & 目标寿命内仍满足功能裕量 & 加宽/加过孔、降额、改变封装 \\\tablerowrule
产品验证 & HTOL、温度循环、功率循环与失效分析 & 样机在目标环境与任务周期下通过 & 修正设计并重新验证受影响范围 \\\tablerowrule
现场运行 & 温压历史、错误趋势、校准与重训计数 & 继续运行、降级、隔离或维护 & 保存证据，更换后用已知负载验收 \\
\bottomrule
\end{longtable}

可靠性的完成边界因此不是“通过一次出厂测试”。设计要先证明额定条件下的寿命模型成立，平台运行时再证明没有长期越过该条件，出现趋势后完成隔离和维护，最后用失效分析把现场证据反馈到下一版设计。这样，器件物理、平台控制和 RAS 才构成完整生命周期。
